\documentclass[11pt]{article}

\usepackage[utf8]{inputenc}
\usepackage[T1]{fontenc}
\usepackage[margin=1in]{geometry}
\usepackage{graphicx}
\usepackage{booktabs}
\usepackage{amsmath}
\usepackage{amssymb}
\usepackage{siunitx}
\usepackage{caption}
\usepackage[hidelinks]{hyperref}

\title{\textbf{Kolmogorov--Arnold Networks for Higgs Boson Classification}\\
\large A comparison with gradient-boosted trees on the FAIR Universe dataset}

\author{Said Abolhassan Razavi\\
\small University of Paris-Saclay --- Internship Research}

\date{June 2026 \\[2pt] \small \textbf{First version} (preliminary report; final version to follow at the end of the internship)}

\begin{document}
\maketitle

\begin{abstract}
\noindent We study Kolmogorov--Arnold Networks (KAN) for the classification of Higgs boson
events ($H\rightarrow\tau\tau$) against background on the full FAIR Universe HiggsML
dataset (2{,}000{,}000 events), benchmarked against XGBoost. Models are evaluated with the
weighted AUC and, as the primary physics metric, the Approximate Median Significance (AMS)
normalised to the full $10\,\mathrm{fb}^{-1}$ luminosity. A properly tuned deep KAN
outperforms XGBoost on both AUC ($0.882$ vs.\ $0.853$) and AMS ($3.69$ vs.\ $2.27$).
An adversarially decorrelated KAN is about $9\times$ more robust to energy-scale
systematics while retaining its accuracy advantage. We additionally explain a known
small-sample artefact in which a strongly regularised XGBoost can appear inferior to KAN.
This is a preliminary report; further analyses are deferred to the final version.
\end{abstract}

\section{Introduction}
The discovery of the Higgs boson in 2012 opened a programme of precision measurements of its
properties. One important channel is the decay into two tau leptons,
$H\rightarrow\tau^+\tau^-$, whose experimental signature is closely mimicked by far more
abundant background processes (notably $Z\rightarrow\tau\tau$, $t\bar t$ and diboson
production). Separating this rare signal from background is a binary classification problem
characterised by a strong class imbalance, where the relevant figure of merit is the
\emph{statistical significance} of the extracted signal rather than raw accuracy.

Gradient-boosted decision trees (XGBoost) are the de facto standard for such tabular
high-energy-physics (HEP) data, and they are typically hard to beat with neural networks.
Kolmogorov--Arnold Networks (KAN) are a recently proposed neural architecture that places
learnable univariate functions on network edges, offering high expressivity and, importantly,
interpretability --- a property of growing interest in physics analyses, where understanding
\emph{why} a model makes a decision matters as much as the decision itself. A further practical
concern in HEP is robustness to systematic uncertainties: a classifier whose output drifts with
detector calibration can bias a measurement.

The objective of this internship is to implement KAN from scratch, evaluate it fairly against
XGBoost on the FAIR Universe HiggsML dataset, and assess its accuracy (AUC), physics
significance (AMS), robustness to systematic uncertainties, and interpretability. We find that,
once properly tuned, KAN meets all of these criteria and surpasses the XGBoost baseline.

\section{Related Work}
KANs were introduced by Liu et al.~\cite{liu2024kan,liu2024kan2} as an alternative to
multilayer perceptrons, replacing fixed node activations with learnable spline functions on
the edges. For tabular data, boosted-tree ensembles such as XGBoost~\cite{chen2016xgboost}
remain very strong baselines. The FAIR Universe HiggsML challenge~\cite{fairuniverse} provides
a large simulated dataset designed to study robustness to systematic uncertainties; the AMS
significance used here follows the asymptotic formulae of Cowan et al.~\cite{cowan2011}.
Robustness to nuisance parameters is addressed via adversarial ``pivoting'' as proposed by
Louppe et al.~\cite{louppe2017}. Preliminary explorations of KANs on this dataset motivated the
present, full-scale and carefully tuned comparison.

\section{Methods}
\paragraph{Dataset.} We use the FAIR Universe ``blackSwan'' dataset (2{,}000{,}000 simulated
$pp$ collisions, 28 features each), split in a stratified manner into 1.4M training, 300k
validation and 300k test events. Missing values are imputed by feature means and inputs are
standardised on the training set.

\paragraph{Metrics.} We report the weighted AUC and, as the primary metric, the AMS,
\begin{equation}
\mathrm{AMS}=\sqrt{2\left[(s+b)\,\ln\!\left(1+\tfrac{s}{b}\right)-s\right]},
\end{equation}
where $s$ and $b$ are the summed signal and background weights passing a decision threshold,
maximised over the threshold. Event weights are normalised so that the test set reproduces the
expected counts of a full $10\,\mathrm{fb}^{-1}$ run ($s=1015$ signal, $b=1{,}050{,}370$
background events, following the dataset's ``LHC Events'' reference).

\paragraph{Models.} The XGBoost baseline uses $8587$ trees of depth $2$ at a low learning rate.
The KAN is implemented from scratch in PyTorch: each edge function is a B-spline evaluated via
the Cox--de Boor recursion, with adaptive knot placement, optional coarse-to-fine
\emph{grid extension} ($G{=}3\rightarrow G{=}10$) and $L_1$ regularisation of the spline
coefficients. The main model is a deep KAN with three layers of width $96$, grid size $G{=}5$,
spline order $k{=}3$, trained with Adam, a cosine learning-rate schedule, mixed precision and
early stopping. An adversarial variant adds a decorrelation network to suppress sensitivity to
the energy scale.

\paragraph{Robustness protocol.} To emulate a mis-calibrated detector we multiply the
momentum-related features by a factor $\gamma\in[0.8,1.2]$ and report the maximal AUC drop
from the nominal $\gamma=1$.

\section{Results and Discussion}
Table~\ref{tab:results} summarises the main results on the held-out test set
(mean $\pm$ standard deviation over 200 bootstrap resamples).

\begin{table}[h]
\centering
\caption{Performance on the full 2M-event test set. AMS is at full LHC luminosity;
``Max drop'' is the AUC degradation under energy-scale shifts (lower is more robust).}
\label{tab:results}
\begin{tabular}{lccc}
\toprule
\textbf{Model} & \textbf{AUC} & \textbf{AMS} & \textbf{Max AUC drop} \\
\midrule
KAN (deep, $[96,96,96]$)      & \textbf{0.8820} & 3.695 & 0.101 \\
KAN $+$ grid extension        & 0.8799 & \textbf{3.935} & 0.090 \\
KAN-Adversarial               & 0.8743 & 2.949 & \textbf{0.006} \\
XGBoost (baseline)            & 0.8533 & 2.269 & 0.052 \\
\bottomrule
\end{tabular}
\end{table}

\paragraph{KAN outperforms XGBoost.} Once properly configured, the deep KAN exceeds XGBoost on
both AUC ($0.882$ vs.\ $0.853$) and the physics-primary AMS ($3.69$ vs.\ $2.27$). The
comparison is fair: XGBoost uses its full, standard configuration on the full dataset. The
gain came primarily from architecture (depth, cubic splines and a higher learning rate); a
shallow, low-learning-rate KAN underperforms. The ROC curves (Fig.~\ref{fig:roc}) show the KAN
lying above XGBoost across the operating range.

\begin{figure}[h]
\centering
\includegraphics[width=0.92\textwidth]{slide_assets/roc_loss.png}
\caption{Left: KAN training loss converging. Right: ROC curves --- the tuned KAN (AUC $0.882$)
lies above XGBoost (AUC $0.853$).}
\label{fig:roc}
\end{figure}

\paragraph{Small-sample artefact.} At very small training sizes ($N\!\lesssim\!2500$) XGBoost
collapses to $\mathrm{AUC}=0.5$. This is not a property of the model but of its default
\texttt{min\_child\_weight}$=1$: because signal weights are tiny ($\sim$$10^{-3}$ per event),
the total signal weight in such a subset falls below this threshold and no split is permitted.
Lowering the threshold restores normal behaviour. This explains why naive small-sample
comparisons can appear to favour KAN, distinct from the genuine, tuned advantage reported here.

\paragraph{Robustness.} The high capacity of the deep KAN makes it more sensitive to the
energy-scale systematic than XGBoost (Fig.~\ref{fig:rob}). Adversarial decorrelation resolves
this: the adversarial KAN reduces the maximal AUC drop to $0.006$ ($\approx9\times$ more robust
than XGBoost) while still beating XGBoost on AUC and AMS, making it the best all-round model.

\begin{figure}[h]
\centering
\includegraphics[width=0.82\textwidth]{slide_assets/robustness.png}
\caption{AUC under energy-scale shifts $\gamma\in[0.8,1.2]$. The deep KAN is most accurate at
$\gamma=1$ but more sensitive away from it.}
\label{fig:rob}
\end{figure}

\paragraph{Interpretability.} Because each KAN edge is an explicit function, feature importance
and learned activations are directly inspectable (Fig.~\ref{fig:feat}). The most important
feature is the visible di-tau invariant mass (\texttt{DER\_mass\_vis}), consistent with
established HEP intuition.

\begin{figure}[h]
\centering
\includegraphics[width=0.82\textwidth]{slide_assets/feature_importance.png}
\caption{Feature importance for KAN and XGBoost. KAN's leading feature, \texttt{DER\_mass\_vis},
is physically meaningful.}
\label{fig:feat}
\end{figure}

\section{Conclusion}
A properly tuned KAN is a competitive and interpretable alternative to XGBoost for
$H\rightarrow\tau\tau$ classification, outperforming it on both AUC and AMS on the full
dataset, with an adversarial variant offering substantially improved robustness to systematic
uncertainties. The central practical finding is that architecture --- network depth, spline
order and learning rate --- determines whether KAN matches or surpasses a strong tree baseline;
the same dataset and a fixed, fairly-configured XGBoost were used throughout.

As this is a preliminary report, the following analyses are planned for the final version:
\begin{itemize}
  \item a systematic hyperparameter and ablation study (depth, grid size, spline order, learning rate);
  \item $k$-fold cross-validation and fuller uncertainty quantification of all metrics;
  \item a more detailed analysis of the robustness--accuracy trade-off, including the
        adversarial strength;
  \item additional baselines and a deeper physics interpretation of the learned spline functions.
\end{itemize}

\begin{thebibliography}{9}
\bibitem{liu2024kan} Z.~Liu et al., \emph{KAN: Kolmogorov--Arnold Networks}, arXiv:2404.19756, 2024.
\bibitem{liu2024kan2} Z.~Liu et al., \emph{KAN 2.0: Kolmogorov--Arnold Networks Meet Science}, arXiv:2408.10205, 2024.
\bibitem{chen2016xgboost} T.~Chen and C.~Guestrin, \emph{XGBoost: A Scalable Tree Boosting System}, KDD, 2016.
\bibitem{cowan2011} G.~Cowan et al., \emph{Asymptotic formulae for likelihood-based tests of new physics}, Eur.\ Phys.\ J.\ C, 2011; arXiv:1007.1727.
\bibitem{louppe2017} G.~Louppe, M.~Kagan and K.~Cranmer, \emph{Learning to Pivot with Adversarial Networks}, NeurIPS, 2017.
\bibitem{fairuniverse} FAIR Universe HiggsML Uncertainty Challenge Dataset, Zenodo:10.5281/zenodo.15131565, 2024.
\end{thebibliography}

\end{document}
